\documentclass[aspectratio=169,10pt]{beamer}
\makeatletter\def\input@path{{../../estilo/}{../estilo/}{estilo/}}\makeatother
\usepackage{devmobile}
\modulo[Flutter]{20}{Tópicos Avançados e Projeto Final}{IA no aparelho e na nuvem, segurança, monetização, outras superfícies, carreira e o projeto integrador}

\begin{document}
\capa

\begin{frame}{Objetivos do módulo}
  \begin{conceito}[Ao final deste módulo você será capaz de]
    \begin{itemize}
      \item Explicar treinamento, inferência e o funcionamento dos modelos de linguagem.
      \item Usar IA no app Flutter: ML Kit, modelos próprios com LiteRT e Gemini via Firebase.
      \item Analisar ameaças com STRIDE e revisar o app com o OWASP MASVS.
      \item Escolher e implementar um modelo de monetização.
      \item Planejar, construir e apresentar o projeto final do curso.
    \end{itemize}
  \end{conceito}
\end{frame}

\sumario

% =============================================================================
%  Fundamentos compartilhados --- Módulo 20: IA, segurança e negócio de apps
%  Incluído pelas trilhas Kotlin e Flutter com % =============================================================================
%  Fundamentos compartilhados --- Módulo 20: IA, segurança e negócio de apps
%  Incluído pelas trilhas Kotlin e Flutter com % =============================================================================
%  Fundamentos compartilhados --- Módulo 20: IA, segurança e negócio de apps
%  Incluído pelas trilhas Kotlin e Flutter com \input{../fundamentos/f20-avancados}
% =============================================================================
\section{Fundamentos de inteligência artificial}

\begin{frame}{IA, aprendizado de máquina e aprendizado profundo}
  \begin{columns}[c]
    \begin{column}{0.44\textwidth}
      \centering
      \begin{tikzpicture}
        \draw[fill=cinzaclaro,draw=azulpetroleo] (0,0) ellipse (2.6 and 1.75);
        \draw[fill=azulpetroleo!25,draw=azulpetroleo] (0.35,-0.25) ellipse (1.95 and 1.3);
        \draw[fill=azulpetroleo!55,draw=azulpetroleo] (0.7,-0.45) ellipse (1.3 and 0.85);
        \draw[fill=ambar,draw=ambar!70!black] (1.0,-0.6) ellipse (0.72 and 0.42);
        \node[font=\scriptsize\bfseries,text=azulpetroleo] at (0,1.42) {IA};
        \node[font=\scriptsize\bfseries,text=azulpetroleo] at (0.35,0.75) {Aprendizado de máquina};
        \node[font=\tiny\bfseries,text=white] at (0.7,0.13) {Aprendizado profundo};
        \node[font=\tiny\bfseries,text=azulpetroleo,align=center] at (1.0,-0.6) {IA\\generativa};
      \end{tikzpicture}
    \end{column}
    \begin{column}{0.54\textwidth}
      \begin{definicao}[Aprendizado de máquina (ML)]
        Em vez de programar as regras, o sistema \textbf{aprende} padrões a partir de exemplos (dados) e os aplica a casos novos.
      \end{definicao}
      \scriptsize
      \begin{tabularx}{\linewidth}{@{}l L@{}}
        \toprule
        \cabecalho{Paradigma} & \cabecalho{Aprende com} \\
        \midrule
        Supervisionado & exemplos rotulados (leitura $\rightarrow$ ``normal'' ou ``falha'') \\
        Não supervisionado & dados sem rótulo, buscando grupos e anomalias \\
        Por reforço & tentativa e erro, guiado por recompensas \\
        \bottomrule
      \end{tabularx}
      \begin{conceito}[Aprendizado profundo]
        ML com redes neurais de muitas camadas; base da visão computacional, da fala e dos modelos de linguagem.
      \end{conceito}
    \end{column}
  \end{columns}
\end{frame}

\begin{frame}{Treinamento e inferência}
  \centering
  \begin{tikzpicture}[node distance=5mm,every node/.style={font=\scriptsize,align=center,minimum height=10mm}]
    \node[caixa,text width=1.7cm] (d) {Dados\\ \tiny rotulados};
    \node[caixapetroleo,right=of d,text width=2.1cm] (t) {Treinamento\\ \tiny servidor com GPU};
    \node[caixaambar,right=of t,text width=1.7cm] (m) {Modelo\\ \tiny pesos};
    \node[caixa,right=of m,text width=2.1cm] (q) {Conversão e quantização};
    \node[caixaescura,right=of q,text width=2.1cm] (i) {Inferência\\ \tiny no app ou na nuvem};
    \foreach \a/\b in {d/t,t/m,m/q,q/i} \draw[seta] (\a) -- (\b);
  \end{tikzpicture}
  \vspace{4pt}
  \begin{columns}[T]
    \begin{column}{0.33\textwidth}
      \begin{definicao}[Treinamento]
        Ajustar milhões de parâmetros (\textbf{pesos}) para que o modelo erre o mínimo nos exemplos. Caro e feito uma vez (ou de tempos em tempos).
      \end{definicao}
    \end{column}
    \begin{column}{0.31\textwidth}
      \begin{definicao}[Inferência]
        Usar o modelo pronto para responder a uma entrada nova. É o que o app faz, milhares de vezes.
      \end{definicao}
    \end{column}
    \begin{column}{0.32\textwidth}
      \begin{pordentro}[Quantização]
        Trocar pesos de 32 bits por inteiros de 8 bits: modelo cerca de 4 vezes menor e mais rápido, com pequena perda de precisão.
      \end{pordentro}
    \end{column}
  \end{columns}
  \begin{atencao}[Avaliar antes de confiar]
    Em detecção de falhas, \textbf{precisão} (dos alertas, quantos eram reais) e \textbf{revocação} (das falhas reais, quantas foram alertadas) importam mais que a acurácia, porque falhas são raras.
  \end{atencao}
\end{frame}

\begin{frame}{Modelos de linguagem (LLMs)}
  \begin{columns}[T]
    \begin{column}{0.5\textwidth}
      \begin{definicao}[LLM --- \emph{Large Language Model}]
        Rede neural treinada com enormes volumes de texto para prever o próximo \textbf{\emph{token}} (pedaço de palavra). Gerando um \emph{token} após o outro, escreve respostas, resume, traduz e segue instruções.
      \end{definicao}
      \scriptsize
      \begin{tabularx}{\linewidth}{@{}l L@{}}
        \toprule
        \cabecalho{Termo} & \cabecalho{Significado} \\
        \midrule
        Janela de contexto & quanto texto o modelo considera de uma vez \\
        Temperatura & grau de aleatoriedade da resposta \\
        Alucinação & resposta fluente, mas falsa \\
        RAG & buscar documentos relevantes e incluí-los na pergunta \\
        Chamada de ferramentas & o modelo pede que o app execute uma função \\
        \bottomrule
      \end{tabularx}
    \end{column}
    \begin{column}{0.48\textwidth}
      \begin{analogia}
        O corretor do teclado levado ao extremo: em vez da próxima palavra, continua parágrafos inteiros.
      \end{analogia}
      \begin{atencao}
        O modelo não ``sabe'' se está certo: produz o texto mais provável. A resposta é uma \textbf{sugestão} que uma pessoa confere.
      \end{atencao}
      \begin{dica}
        Com RAG, o LabMonitor responde com base nos manuais dos equipamentos, citando a fonte.
      \end{dica}
    \end{column}
  \end{columns}
\end{frame}

\begin{frame}[fragile]{Engenharia de \emph{prompt} e uso responsável}
  \begin{columns}[T]
    \begin{column}{0.52\textwidth}
\begin{codigo}[title=Estrutura de um bom \emph{prompt}]{bash}
Papel:    Você é assistente técnico de um
          laboratório de eletroeletrônica.
Contexto: Leituras da bancada B3 (última hora):
          {leituras}
Tarefa:   Explique o alerta de temperatura.
Formato:  No máximo 3 frases, em português,
          e termine com uma ação recomendada.
Limites:  Se os dados não bastarem, diga isso.
\end{codigo}
      \begin{dica}
        Dê exemplos de entrada e saída esperadas (\emph{few-shot}) e peça respostas em JSON quando o app precisar processá-las.
      \end{dica}
    \end{column}
    \begin{column}{0.46\textwidth}
      \begin{erro}[Injeção de \emph{prompt}]
        Texto vindo do usuário ou de documentos pode conter instruções (``ignore as regras anteriores\ldots''). Nunca dê ao modelo poder de executar ações sensíveis sem confirmação e validação no servidor.
      \end{erro}
      \begin{conceito}[Uso responsável]
        \begin{itemize}
          \item Informe que a resposta foi gerada por IA.
          \item Minimize dados pessoais enviados (LGPD).
          \item Limite custos: cada \emph{token} é cobrado.
          \item Avalie com casos de teste, como qualquer código.
        \end{itemize}
      \end{conceito}
    \end{column}
  \end{columns}
\end{frame}

% =============================================================================
\section{Fundamentos de segurança da informação}

\begin{frame}{Princípios de segurança}
  \begin{columns}[c]
    \begin{column}{0.42\textwidth}
      \centering
      \begin{tikzpicture}
        \draw[fill=azulpetroleo!8,draw=azulpetroleo,line width=1pt] (90:1.7) -- (210:1.7) -- (330:1.7) -- cycle;
        \node[caixapetroleo,font=\scriptsize] at (90:1.95) {Confidencialidade};
        \node[caixaambar,font=\scriptsize] at (210:2.0) {Integridade};
        \node[caixaescura,font=\scriptsize] at (330:2.0) {Disponibilidade};
        \node[font=\scriptsize\bfseries,text=azulpetroleo] at (0,0) {CIA};
      \end{tikzpicture}
      \\[2pt]{\tiny Tríade CID (em inglês, CIA).}
    \end{column}
    \begin{column}{0.56\textwidth}
      \scriptsize
      \begin{tabularx}{\linewidth}{@{}l L@{}}
        \toprule
        \cabecalho{Propriedade} & \cabecalho{Pergunta} \\
        \midrule
        Confidencialidade & só quem pode vê os dados? \\
        Integridade & os dados não foram alterados indevidamente? \\
        Disponibilidade & o serviço funciona quando precisam dele? \\
        Autenticidade & quem diz ser é quem é? \\
        \bottomrule
      \end{tabularx}
      \begin{conceito}[Princípios de projeto]
        \textbf{Menor privilégio} (cada parte com o mínimo de acesso), \textbf{defesa em profundidade} (várias camadas), \textbf{seguro por padrão} e \textbf{nunca confiar no cliente}: o app pode ser modificado por quem o instala.
      \end{conceito}
    \end{column}
  \end{columns}
\end{frame}

\begin{frame}{Criptografia essencial}
  \footnotesize
  \begin{tabularx}{\textwidth}{@{}p{2.5cm} L L L@{}}
    \toprule
    & \cabecalho{Hash} & \cabecalho{Simétrica} & \cabecalho{Assimétrica} \\
    \midrule
    Ideia & resumo de tamanho fixo, impossível de reverter & a mesma chave cifra e decifra & par de chaves: pública e privada \\
    Exemplos & SHA-256; para senhas, Argon2 ou bcrypt & AES-GCM & RSA, curvas elípticas (ECDSA, X25519) \\
    Uso no app & integridade de arquivos, armazenar senhas no servidor & dados locais cifrados, conteúdo do canal TLS & assinatura do app, troca de chaves, \emph{passkeys} \\
    \bottomrule
  \end{tabularx}
  \vspace{4pt}
  \begin{columns}[T]
    \begin{column}{0.48\textwidth}
      \begin{pordentro}[Como o HTTPS junta tudo]
        O TLS usa criptografia assimétrica para autenticar o servidor (certificado) e combinar uma chave secreta; depois cifra o tráfego com criptografia simétrica, que é muito mais rápida.
      \end{pordentro}
    \end{column}
    \begin{column}{0.48\textwidth}
      \begin{erro}[Criptografia caseira]
        Inventar algoritmos, reutilizar vetores de inicialização ou guardar a chave junto com o dado cifrado. Use bibliotecas e o cofre de chaves do sistema (Keystore, Keychain).
      \end{erro}
    \end{column}
  \end{columns}
\end{frame}

\begin{frame}{Modelagem de ameaças com STRIDE}
  \begin{columns}[T]
    \begin{column}{0.62\textwidth}
      \scriptsize
      \begin{tabularx}{\linewidth}{@{}l L L@{}}
        \toprule
        \cabecalho{Ameaça} & \cabecalho{Exemplo no LabMonitor} & \cabecalho{Defesa} \\
        \midrule
        \textbf{S}poofing (falsificação) & alguém se passa por um ESP32 e envia leituras falsas & autenticar dispositivos (certificados, MQTT com senha) \\
        \textbf{T}ampering (adulteração) & alterar a reserva no tráfego & TLS, validação no servidor \\
        \textbf{R}epudiation (repúdio) & aluno nega ter cancelado a reserva & registros de auditoria \\
        \textbf{I}nformation disclosure & token salvo em texto puro & armazenamento seguro \\
        \textbf{D}enial of service & app em laço enviando milhares de pedidos & limites de taxa no servidor \\
        \textbf{E}levation of privilege & aluno vira administrador alterando o app & regras e papéis no servidor \\
        \bottomrule
      \end{tabularx}
    \end{column}
    \begin{column}{0.36\textwidth}
      \begin{definicao}[Modelagem de ameaças]
        Analisar sistematicamente o que pode dar errado, \textbf{antes} de construir: desenhar o fluxo de dados, listar ameaças, priorizar e definir defesas.
      \end{definicao}
      \begin{historia}
        STRIDE foi criado na Microsoft no fim dos anos 1990 e virou uma das técnicas mais ensinadas.
      \end{historia}
    \end{column}
  \end{columns}
\end{frame}

\begin{frame}{OWASP Mobile Top 10 e MASVS}
  \begin{columns}[T]
    \begin{column}{0.5\textwidth}
      \scriptsize
      \begin{tabularx}{\linewidth}{@{}l L@{}}
        \toprule
        \multicolumn{2}{@{}l}{\cabecalho{OWASP Mobile Top 10 (2024)}} \\
        \midrule
        M1 & uso inadequado de credenciais \\
        M2 & segurança insuficiente da cadeia de suprimentos \\
        M3 & autenticação e autorização inseguras \\
        M4 & validação insuficiente de entradas e saídas \\
        M5 & comunicação insegura \\
        M6 & controles de privacidade inadequados \\
        M7 & proteção insuficiente do binário \\
        M8 & configuração de segurança incorreta \\
        M9 & armazenamento de dados inseguro \\
        M10 & criptografia insuficiente \\
        \bottomrule
      \end{tabularx}
    \end{column}
    \begin{column}{0.48\textwidth}
      \begin{definicao}[OWASP]
        Fundação aberta que publica guias de segurança. O \textbf{Top 10} lista os riscos mais comuns; o \textbf{MASVS} define os requisitos de segurança para apps, e o \textbf{MASTG} ensina a testá-los.
      \end{definicao}
      \begin{conceito}[Grupos do MASVS]
        Armazenamento, criptografia, autenticação, rede, plataforma, código, resiliência e privacidade.
      \end{conceito}
      \begin{dica}
        Use o MASVS como \emph{checklist} de revisão antes de cada versão importante.
      \end{dica}
    \end{column}
  \end{columns}
\end{frame}

% =============================================================================
\section{Fundamentos de produto e negócio}

\begin{frame}{Do problema ao MVP}
  \begin{columns}[c]
    \begin{column}{0.4\textwidth}
      \centering
      \begin{tikzpicture}
        \node[circle,fill=azulpetroleo,text=white,font=\scriptsize\bfseries,minimum size=1.3cm] (c) at (90:1.3) {Construir};
        \node[circle,fill=ambar,text=azulpetroleo,font=\scriptsize\bfseries,minimum size=1.3cm] (m) at (-30:1.3) {Medir};
        \node[circle,fill=azulprofundo,text=white,font=\scriptsize\bfseries,minimum size=1.3cm] (a) at (210:1.3) {Aprender};
        \draw[setaambar] (c) to[bend left=25] node[rotulo,right,font=\tiny]{produto} (m);
        \draw[setaambar] (m) to[bend left=25] node[rotulo,below,font=\tiny]{dados} (a);
        \draw[setaambar] (a) to[bend left=25] node[rotulo,left,font=\tiny]{ideias} (c);
      \end{tikzpicture}
    \end{column}
    \begin{column}{0.58\textwidth}
      \begin{definicao}[MVP --- Produto Mínimo Viável]
        A menor versão do produto que permite \textbf{aprender} com usuários reais se a solução resolve o problema. Não é um produto mal feito: é um produto pequeno.
      \end{definicao}
      \begin{conceito}[Métricas AARRR]
        \textbf{A}quisição (chegam), \textbf{A}tivação (têm uma boa primeira experiência), \textbf{R}etenção (voltam), \textbf{R}eceita (pagam) e \textbf{R}ecomendação (indicam).
      \end{conceito}
      \begin{historia}
        O ciclo construir--medir--aprender foi popularizado por Eric Ries no livro \emph{A Startup Enxuta} (2011).
      \end{historia}
    \end{column}
  \end{columns}
\end{frame}

\begin{frame}{Modelos de receita e métricas do negócio}
  \begin{columns}[T]
    \begin{column}{0.46\textwidth}
      \scriptsize
      \begin{tabularx}{\linewidth}{@{}l L@{}}
        \toprule
        \cabecalho{Modelo} & \cabecalho{Como gera receita} \\
        \midrule
        Pago & paga-se para instalar \\
        Freemium & grátis, com recursos pagos \\
        Assinatura & cobrança recorrente \\
        Anúncios & exibições e cliques \\
        B2B / institucional & a organização paga pelo serviço \\
        Público / acadêmico & financiado por projeto ou órgão \\
        \bottomrule
      \end{tabularx}
    \end{column}
    \begin{column}{0.52\textwidth}
      \scriptsize
      \begin{tabularx}{\linewidth}{@{}l L@{}}
        \toprule
        \cabecalho{Métrica} & \cabecalho{Definição} \\
        \midrule
        ARPU & receita média por usuário ativo no período \\
        \emph{Churn} & \% de assinantes que cancelam por mês \\
        LTV & receita esperada de um usuário durante toda a relação; aproximação: ARPU mensal $\times$ margem $\div$ \emph{churn} mensal \\
        CAC & custo para adquirir um usuário (marketing $\div$ novos usuários) \\
        \bottomrule
      \end{tabularx}
      \begin{exemplo}[Conta rápida]
        ARPU de R\$\,10/mês, margem de 70\,\% e \emph{churn} de 5\,\%/mês: LTV $\approx 10 \times 0{,}7 \div 0{,}05 =$ R\$\,140. Gastar R\$\,200 para adquirir cada usuário dá prejuízo.
      \end{exemplo}
    \end{column}
  \end{columns}
\end{frame}

% =============================================================================
\section{Fundamentos de inteligência artificial}

\begin{frame}{IA, aprendizado de máquina e aprendizado profundo}
  \begin{columns}[c]
    \begin{column}{0.44\textwidth}
      \centering
      \begin{tikzpicture}
        \draw[fill=cinzaclaro,draw=azulpetroleo] (0,0) ellipse (2.6 and 1.75);
        \draw[fill=azulpetroleo!25,draw=azulpetroleo] (0.35,-0.25) ellipse (1.95 and 1.3);
        \draw[fill=azulpetroleo!55,draw=azulpetroleo] (0.7,-0.45) ellipse (1.3 and 0.85);
        \draw[fill=ambar,draw=ambar!70!black] (1.0,-0.6) ellipse (0.72 and 0.42);
        \node[font=\scriptsize\bfseries,text=azulpetroleo] at (0,1.42) {IA};
        \node[font=\scriptsize\bfseries,text=azulpetroleo] at (0.35,0.75) {Aprendizado de máquina};
        \node[font=\tiny\bfseries,text=white] at (0.7,0.13) {Aprendizado profundo};
        \node[font=\tiny\bfseries,text=azulpetroleo,align=center] at (1.0,-0.6) {IA\\generativa};
      \end{tikzpicture}
    \end{column}
    \begin{column}{0.54\textwidth}
      \begin{definicao}[Aprendizado de máquina (ML)]
        Em vez de programar as regras, o sistema \textbf{aprende} padrões a partir de exemplos (dados) e os aplica a casos novos.
      \end{definicao}
      \scriptsize
      \begin{tabularx}{\linewidth}{@{}l L@{}}
        \toprule
        \cabecalho{Paradigma} & \cabecalho{Aprende com} \\
        \midrule
        Supervisionado & exemplos rotulados (leitura $\rightarrow$ ``normal'' ou ``falha'') \\
        Não supervisionado & dados sem rótulo, buscando grupos e anomalias \\
        Por reforço & tentativa e erro, guiado por recompensas \\
        \bottomrule
      \end{tabularx}
      \begin{conceito}[Aprendizado profundo]
        ML com redes neurais de muitas camadas; base da visão computacional, da fala e dos modelos de linguagem.
      \end{conceito}
    \end{column}
  \end{columns}
\end{frame}

\begin{frame}{Treinamento e inferência}
  \centering
  \begin{tikzpicture}[node distance=5mm,every node/.style={font=\scriptsize,align=center,minimum height=10mm}]
    \node[caixa,text width=1.7cm] (d) {Dados\\ \tiny rotulados};
    \node[caixapetroleo,right=of d,text width=2.1cm] (t) {Treinamento\\ \tiny servidor com GPU};
    \node[caixaambar,right=of t,text width=1.7cm] (m) {Modelo\\ \tiny pesos};
    \node[caixa,right=of m,text width=2.1cm] (q) {Conversão e quantização};
    \node[caixaescura,right=of q,text width=2.1cm] (i) {Inferência\\ \tiny no app ou na nuvem};
    \foreach \a/\b in {d/t,t/m,m/q,q/i} \draw[seta] (\a) -- (\b);
  \end{tikzpicture}
  \vspace{4pt}
  \begin{columns}[T]
    \begin{column}{0.33\textwidth}
      \begin{definicao}[Treinamento]
        Ajustar milhões de parâmetros (\textbf{pesos}) para que o modelo erre o mínimo nos exemplos. Caro e feito uma vez (ou de tempos em tempos).
      \end{definicao}
    \end{column}
    \begin{column}{0.31\textwidth}
      \begin{definicao}[Inferência]
        Usar o modelo pronto para responder a uma entrada nova. É o que o app faz, milhares de vezes.
      \end{definicao}
    \end{column}
    \begin{column}{0.32\textwidth}
      \begin{pordentro}[Quantização]
        Trocar pesos de 32 bits por inteiros de 8 bits: modelo cerca de 4 vezes menor e mais rápido, com pequena perda de precisão.
      \end{pordentro}
    \end{column}
  \end{columns}
  \begin{atencao}[Avaliar antes de confiar]
    Em detecção de falhas, \textbf{precisão} (dos alertas, quantos eram reais) e \textbf{revocação} (das falhas reais, quantas foram alertadas) importam mais que a acurácia, porque falhas são raras.
  \end{atencao}
\end{frame}

\begin{frame}{Modelos de linguagem (LLMs)}
  \begin{columns}[T]
    \begin{column}{0.5\textwidth}
      \begin{definicao}[LLM --- \emph{Large Language Model}]
        Rede neural treinada com enormes volumes de texto para prever o próximo \textbf{\emph{token}} (pedaço de palavra). Gerando um \emph{token} após o outro, escreve respostas, resume, traduz e segue instruções.
      \end{definicao}
      \scriptsize
      \begin{tabularx}{\linewidth}{@{}l L@{}}
        \toprule
        \cabecalho{Termo} & \cabecalho{Significado} \\
        \midrule
        Janela de contexto & quanto texto o modelo considera de uma vez \\
        Temperatura & grau de aleatoriedade da resposta \\
        Alucinação & resposta fluente, mas falsa \\
        RAG & buscar documentos relevantes e incluí-los na pergunta \\
        Chamada de ferramentas & o modelo pede que o app execute uma função \\
        \bottomrule
      \end{tabularx}
    \end{column}
    \begin{column}{0.48\textwidth}
      \begin{analogia}
        O corretor do teclado levado ao extremo: em vez da próxima palavra, continua parágrafos inteiros.
      \end{analogia}
      \begin{atencao}
        O modelo não ``sabe'' se está certo: produz o texto mais provável. A resposta é uma \textbf{sugestão} que uma pessoa confere.
      \end{atencao}
      \begin{dica}
        Com RAG, o LabMonitor responde com base nos manuais dos equipamentos, citando a fonte.
      \end{dica}
    \end{column}
  \end{columns}
\end{frame}

\begin{frame}[fragile]{Engenharia de \emph{prompt} e uso responsável}
  \begin{columns}[T]
    \begin{column}{0.52\textwidth}
\begin{codigo}[title=Estrutura de um bom \emph{prompt}]{bash}
Papel:    Você é assistente técnico de um
          laboratório de eletroeletrônica.
Contexto: Leituras da bancada B3 (última hora):
          {leituras}
Tarefa:   Explique o alerta de temperatura.
Formato:  No máximo 3 frases, em português,
          e termine com uma ação recomendada.
Limites:  Se os dados não bastarem, diga isso.
\end{codigo}
      \begin{dica}
        Dê exemplos de entrada e saída esperadas (\emph{few-shot}) e peça respostas em JSON quando o app precisar processá-las.
      \end{dica}
    \end{column}
    \begin{column}{0.46\textwidth}
      \begin{erro}[Injeção de \emph{prompt}]
        Texto vindo do usuário ou de documentos pode conter instruções (``ignore as regras anteriores\ldots''). Nunca dê ao modelo poder de executar ações sensíveis sem confirmação e validação no servidor.
      \end{erro}
      \begin{conceito}[Uso responsável]
        \begin{itemize}
          \item Informe que a resposta foi gerada por IA.
          \item Minimize dados pessoais enviados (LGPD).
          \item Limite custos: cada \emph{token} é cobrado.
          \item Avalie com casos de teste, como qualquer código.
        \end{itemize}
      \end{conceito}
    \end{column}
  \end{columns}
\end{frame}

% =============================================================================
\section{Fundamentos de segurança da informação}

\begin{frame}{Princípios de segurança}
  \begin{columns}[c]
    \begin{column}{0.42\textwidth}
      \centering
      \begin{tikzpicture}
        \draw[fill=azulpetroleo!8,draw=azulpetroleo,line width=1pt] (90:1.7) -- (210:1.7) -- (330:1.7) -- cycle;
        \node[caixapetroleo,font=\scriptsize] at (90:1.95) {Confidencialidade};
        \node[caixaambar,font=\scriptsize] at (210:2.0) {Integridade};
        \node[caixaescura,font=\scriptsize] at (330:2.0) {Disponibilidade};
        \node[font=\scriptsize\bfseries,text=azulpetroleo] at (0,0) {CIA};
      \end{tikzpicture}
      \\[2pt]{\tiny Tríade CID (em inglês, CIA).}
    \end{column}
    \begin{column}{0.56\textwidth}
      \scriptsize
      \begin{tabularx}{\linewidth}{@{}l L@{}}
        \toprule
        \cabecalho{Propriedade} & \cabecalho{Pergunta} \\
        \midrule
        Confidencialidade & só quem pode vê os dados? \\
        Integridade & os dados não foram alterados indevidamente? \\
        Disponibilidade & o serviço funciona quando precisam dele? \\
        Autenticidade & quem diz ser é quem é? \\
        \bottomrule
      \end{tabularx}
      \begin{conceito}[Princípios de projeto]
        \textbf{Menor privilégio} (cada parte com o mínimo de acesso), \textbf{defesa em profundidade} (várias camadas), \textbf{seguro por padrão} e \textbf{nunca confiar no cliente}: o app pode ser modificado por quem o instala.
      \end{conceito}
    \end{column}
  \end{columns}
\end{frame}

\begin{frame}{Criptografia essencial}
  \footnotesize
  \begin{tabularx}{\textwidth}{@{}p{2.5cm} L L L@{}}
    \toprule
    & \cabecalho{Hash} & \cabecalho{Simétrica} & \cabecalho{Assimétrica} \\
    \midrule
    Ideia & resumo de tamanho fixo, impossível de reverter & a mesma chave cifra e decifra & par de chaves: pública e privada \\
    Exemplos & SHA-256; para senhas, Argon2 ou bcrypt & AES-GCM & RSA, curvas elípticas (ECDSA, X25519) \\
    Uso no app & integridade de arquivos, armazenar senhas no servidor & dados locais cifrados, conteúdo do canal TLS & assinatura do app, troca de chaves, \emph{passkeys} \\
    \bottomrule
  \end{tabularx}
  \vspace{4pt}
  \begin{columns}[T]
    \begin{column}{0.48\textwidth}
      \begin{pordentro}[Como o HTTPS junta tudo]
        O TLS usa criptografia assimétrica para autenticar o servidor (certificado) e combinar uma chave secreta; depois cifra o tráfego com criptografia simétrica, que é muito mais rápida.
      \end{pordentro}
    \end{column}
    \begin{column}{0.48\textwidth}
      \begin{erro}[Criptografia caseira]
        Inventar algoritmos, reutilizar vetores de inicialização ou guardar a chave junto com o dado cifrado. Use bibliotecas e o cofre de chaves do sistema (Keystore, Keychain).
      \end{erro}
    \end{column}
  \end{columns}
\end{frame}

\begin{frame}{Modelagem de ameaças com STRIDE}
  \begin{columns}[T]
    \begin{column}{0.62\textwidth}
      \scriptsize
      \begin{tabularx}{\linewidth}{@{}l L L@{}}
        \toprule
        \cabecalho{Ameaça} & \cabecalho{Exemplo no LabMonitor} & \cabecalho{Defesa} \\
        \midrule
        \textbf{S}poofing (falsificação) & alguém se passa por um ESP32 e envia leituras falsas & autenticar dispositivos (certificados, MQTT com senha) \\
        \textbf{T}ampering (adulteração) & alterar a reserva no tráfego & TLS, validação no servidor \\
        \textbf{R}epudiation (repúdio) & aluno nega ter cancelado a reserva & registros de auditoria \\
        \textbf{I}nformation disclosure & token salvo em texto puro & armazenamento seguro \\
        \textbf{D}enial of service & app em laço enviando milhares de pedidos & limites de taxa no servidor \\
        \textbf{E}levation of privilege & aluno vira administrador alterando o app & regras e papéis no servidor \\
        \bottomrule
      \end{tabularx}
    \end{column}
    \begin{column}{0.36\textwidth}
      \begin{definicao}[Modelagem de ameaças]
        Analisar sistematicamente o que pode dar errado, \textbf{antes} de construir: desenhar o fluxo de dados, listar ameaças, priorizar e definir defesas.
      \end{definicao}
      \begin{historia}
        STRIDE foi criado na Microsoft no fim dos anos 1990 e virou uma das técnicas mais ensinadas.
      \end{historia}
    \end{column}
  \end{columns}
\end{frame}

\begin{frame}{OWASP Mobile Top 10 e MASVS}
  \begin{columns}[T]
    \begin{column}{0.5\textwidth}
      \scriptsize
      \begin{tabularx}{\linewidth}{@{}l L@{}}
        \toprule
        \multicolumn{2}{@{}l}{\cabecalho{OWASP Mobile Top 10 (2024)}} \\
        \midrule
        M1 & uso inadequado de credenciais \\
        M2 & segurança insuficiente da cadeia de suprimentos \\
        M3 & autenticação e autorização inseguras \\
        M4 & validação insuficiente de entradas e saídas \\
        M5 & comunicação insegura \\
        M6 & controles de privacidade inadequados \\
        M7 & proteção insuficiente do binário \\
        M8 & configuração de segurança incorreta \\
        M9 & armazenamento de dados inseguro \\
        M10 & criptografia insuficiente \\
        \bottomrule
      \end{tabularx}
    \end{column}
    \begin{column}{0.48\textwidth}
      \begin{definicao}[OWASP]
        Fundação aberta que publica guias de segurança. O \textbf{Top 10} lista os riscos mais comuns; o \textbf{MASVS} define os requisitos de segurança para apps, e o \textbf{MASTG} ensina a testá-los.
      \end{definicao}
      \begin{conceito}[Grupos do MASVS]
        Armazenamento, criptografia, autenticação, rede, plataforma, código, resiliência e privacidade.
      \end{conceito}
      \begin{dica}
        Use o MASVS como \emph{checklist} de revisão antes de cada versão importante.
      \end{dica}
    \end{column}
  \end{columns}
\end{frame}

% =============================================================================
\section{Fundamentos de produto e negócio}

\begin{frame}{Do problema ao MVP}
  \begin{columns}[c]
    \begin{column}{0.4\textwidth}
      \centering
      \begin{tikzpicture}
        \node[circle,fill=azulpetroleo,text=white,font=\scriptsize\bfseries,minimum size=1.3cm] (c) at (90:1.3) {Construir};
        \node[circle,fill=ambar,text=azulpetroleo,font=\scriptsize\bfseries,minimum size=1.3cm] (m) at (-30:1.3) {Medir};
        \node[circle,fill=azulprofundo,text=white,font=\scriptsize\bfseries,minimum size=1.3cm] (a) at (210:1.3) {Aprender};
        \draw[setaambar] (c) to[bend left=25] node[rotulo,right,font=\tiny]{produto} (m);
        \draw[setaambar] (m) to[bend left=25] node[rotulo,below,font=\tiny]{dados} (a);
        \draw[setaambar] (a) to[bend left=25] node[rotulo,left,font=\tiny]{ideias} (c);
      \end{tikzpicture}
    \end{column}
    \begin{column}{0.58\textwidth}
      \begin{definicao}[MVP --- Produto Mínimo Viável]
        A menor versão do produto que permite \textbf{aprender} com usuários reais se a solução resolve o problema. Não é um produto mal feito: é um produto pequeno.
      \end{definicao}
      \begin{conceito}[Métricas AARRR]
        \textbf{A}quisição (chegam), \textbf{A}tivação (têm uma boa primeira experiência), \textbf{R}etenção (voltam), \textbf{R}eceita (pagam) e \textbf{R}ecomendação (indicam).
      \end{conceito}
      \begin{historia}
        O ciclo construir--medir--aprender foi popularizado por Eric Ries no livro \emph{A Startup Enxuta} (2011).
      \end{historia}
    \end{column}
  \end{columns}
\end{frame}

\begin{frame}{Modelos de receita e métricas do negócio}
  \begin{columns}[T]
    \begin{column}{0.46\textwidth}
      \scriptsize
      \begin{tabularx}{\linewidth}{@{}l L@{}}
        \toprule
        \cabecalho{Modelo} & \cabecalho{Como gera receita} \\
        \midrule
        Pago & paga-se para instalar \\
        Freemium & grátis, com recursos pagos \\
        Assinatura & cobrança recorrente \\
        Anúncios & exibições e cliques \\
        B2B / institucional & a organização paga pelo serviço \\
        Público / acadêmico & financiado por projeto ou órgão \\
        \bottomrule
      \end{tabularx}
    \end{column}
    \begin{column}{0.52\textwidth}
      \scriptsize
      \begin{tabularx}{\linewidth}{@{}l L@{}}
        \toprule
        \cabecalho{Métrica} & \cabecalho{Definição} \\
        \midrule
        ARPU & receita média por usuário ativo no período \\
        \emph{Churn} & \% de assinantes que cancelam por mês \\
        LTV & receita esperada de um usuário durante toda a relação; aproximação: ARPU mensal $\times$ margem $\div$ \emph{churn} mensal \\
        CAC & custo para adquirir um usuário (marketing $\div$ novos usuários) \\
        \bottomrule
      \end{tabularx}
      \begin{exemplo}[Conta rápida]
        ARPU de R\$\,10/mês, margem de 70\,\% e \emph{churn} de 5\,\%/mês: LTV $\approx 10 \times 0{,}7 \div 0{,}05 =$ R\$\,140. Gastar R\$\,200 para adquirir cada usuário dá prejuízo.
      \end{exemplo}
    \end{column}
  \end{columns}
\end{frame}

% =============================================================================
\section{Fundamentos de inteligência artificial}

\begin{frame}{IA, aprendizado de máquina e aprendizado profundo}
  \begin{columns}[c]
    \begin{column}{0.44\textwidth}
      \centering
      \begin{tikzpicture}
        \draw[fill=cinzaclaro,draw=azulpetroleo] (0,0) ellipse (2.6 and 1.75);
        \draw[fill=azulpetroleo!25,draw=azulpetroleo] (0.35,-0.25) ellipse (1.95 and 1.3);
        \draw[fill=azulpetroleo!55,draw=azulpetroleo] (0.7,-0.45) ellipse (1.3 and 0.85);
        \draw[fill=ambar,draw=ambar!70!black] (1.0,-0.6) ellipse (0.72 and 0.42);
        \node[font=\scriptsize\bfseries,text=azulpetroleo] at (0,1.42) {IA};
        \node[font=\scriptsize\bfseries,text=azulpetroleo] at (0.35,0.75) {Aprendizado de máquina};
        \node[font=\tiny\bfseries,text=white] at (0.7,0.13) {Aprendizado profundo};
        \node[font=\tiny\bfseries,text=azulpetroleo,align=center] at (1.0,-0.6) {IA\\generativa};
      \end{tikzpicture}
    \end{column}
    \begin{column}{0.54\textwidth}
      \begin{definicao}[Aprendizado de máquina (ML)]
        Em vez de programar as regras, o sistema \textbf{aprende} padrões a partir de exemplos (dados) e os aplica a casos novos.
      \end{definicao}
      \scriptsize
      \begin{tabularx}{\linewidth}{@{}l L@{}}
        \toprule
        \cabecalho{Paradigma} & \cabecalho{Aprende com} \\
        \midrule
        Supervisionado & exemplos rotulados (leitura $\rightarrow$ ``normal'' ou ``falha'') \\
        Não supervisionado & dados sem rótulo, buscando grupos e anomalias \\
        Por reforço & tentativa e erro, guiado por recompensas \\
        \bottomrule
      \end{tabularx}
      \begin{conceito}[Aprendizado profundo]
        ML com redes neurais de muitas camadas; base da visão computacional, da fala e dos modelos de linguagem.
      \end{conceito}
    \end{column}
  \end{columns}
\end{frame}

\begin{frame}{Treinamento e inferência}
  \centering
  \begin{tikzpicture}[node distance=5mm,every node/.style={font=\scriptsize,align=center,minimum height=10mm}]
    \node[caixa,text width=1.7cm] (d) {Dados\\ \tiny rotulados};
    \node[caixapetroleo,right=of d,text width=2.1cm] (t) {Treinamento\\ \tiny servidor com GPU};
    \node[caixaambar,right=of t,text width=1.7cm] (m) {Modelo\\ \tiny pesos};
    \node[caixa,right=of m,text width=2.1cm] (q) {Conversão e quantização};
    \node[caixaescura,right=of q,text width=2.1cm] (i) {Inferência\\ \tiny no app ou na nuvem};
    \foreach \a/\b in {d/t,t/m,m/q,q/i} \draw[seta] (\a) -- (\b);
  \end{tikzpicture}
  \vspace{4pt}
  \begin{columns}[T]
    \begin{column}{0.33\textwidth}
      \begin{definicao}[Treinamento]
        Ajustar milhões de parâmetros (\textbf{pesos}) para que o modelo erre o mínimo nos exemplos. Caro e feito uma vez (ou de tempos em tempos).
      \end{definicao}
    \end{column}
    \begin{column}{0.31\textwidth}
      \begin{definicao}[Inferência]
        Usar o modelo pronto para responder a uma entrada nova. É o que o app faz, milhares de vezes.
      \end{definicao}
    \end{column}
    \begin{column}{0.32\textwidth}
      \begin{pordentro}[Quantização]
        Trocar pesos de 32 bits por inteiros de 8 bits: modelo cerca de 4 vezes menor e mais rápido, com pequena perda de precisão.
      \end{pordentro}
    \end{column}
  \end{columns}
  \begin{atencao}[Avaliar antes de confiar]
    Em detecção de falhas, \textbf{precisão} (dos alertas, quantos eram reais) e \textbf{revocação} (das falhas reais, quantas foram alertadas) importam mais que a acurácia, porque falhas são raras.
  \end{atencao}
\end{frame}

\begin{frame}{Modelos de linguagem (LLMs)}
  \begin{columns}[T]
    \begin{column}{0.5\textwidth}
      \begin{definicao}[LLM --- \emph{Large Language Model}]
        Rede neural treinada com enormes volumes de texto para prever o próximo \textbf{\emph{token}} (pedaço de palavra). Gerando um \emph{token} após o outro, escreve respostas, resume, traduz e segue instruções.
      \end{definicao}
      \scriptsize
      \begin{tabularx}{\linewidth}{@{}l L@{}}
        \toprule
        \cabecalho{Termo} & \cabecalho{Significado} \\
        \midrule
        Janela de contexto & quanto texto o modelo considera de uma vez \\
        Temperatura & grau de aleatoriedade da resposta \\
        Alucinação & resposta fluente, mas falsa \\
        RAG & buscar documentos relevantes e incluí-los na pergunta \\
        Chamada de ferramentas & o modelo pede que o app execute uma função \\
        \bottomrule
      \end{tabularx}
    \end{column}
    \begin{column}{0.48\textwidth}
      \begin{analogia}
        O corretor do teclado levado ao extremo: em vez da próxima palavra, continua parágrafos inteiros.
      \end{analogia}
      \begin{atencao}
        O modelo não ``sabe'' se está certo: produz o texto mais provável. A resposta é uma \textbf{sugestão} que uma pessoa confere.
      \end{atencao}
      \begin{dica}
        Com RAG, o LabMonitor responde com base nos manuais dos equipamentos, citando a fonte.
      \end{dica}
    \end{column}
  \end{columns}
\end{frame}

\begin{frame}[fragile]{Engenharia de \emph{prompt} e uso responsável}
  \begin{columns}[T]
    \begin{column}{0.52\textwidth}
\begin{codigo}[title=Estrutura de um bom \emph{prompt}]{bash}
Papel:    Você é assistente técnico de um
          laboratório de eletroeletrônica.
Contexto: Leituras da bancada B3 (última hora):
          {leituras}
Tarefa:   Explique o alerta de temperatura.
Formato:  No máximo 3 frases, em português,
          e termine com uma ação recomendada.
Limites:  Se os dados não bastarem, diga isso.
\end{codigo}
      \begin{dica}
        Dê exemplos de entrada e saída esperadas (\emph{few-shot}) e peça respostas em JSON quando o app precisar processá-las.
      \end{dica}
    \end{column}
    \begin{column}{0.46\textwidth}
      \begin{erro}[Injeção de \emph{prompt}]
        Texto vindo do usuário ou de documentos pode conter instruções (``ignore as regras anteriores\ldots''). Nunca dê ao modelo poder de executar ações sensíveis sem confirmação e validação no servidor.
      \end{erro}
      \begin{conceito}[Uso responsável]
        \begin{itemize}
          \item Informe que a resposta foi gerada por IA.
          \item Minimize dados pessoais enviados (LGPD).
          \item Limite custos: cada \emph{token} é cobrado.
          \item Avalie com casos de teste, como qualquer código.
        \end{itemize}
      \end{conceito}
    \end{column}
  \end{columns}
\end{frame}

% =============================================================================
\section{Fundamentos de segurança da informação}

\begin{frame}{Princípios de segurança}
  \begin{columns}[c]
    \begin{column}{0.42\textwidth}
      \centering
      \begin{tikzpicture}
        \draw[fill=azulpetroleo!8,draw=azulpetroleo,line width=1pt] (90:1.7) -- (210:1.7) -- (330:1.7) -- cycle;
        \node[caixapetroleo,font=\scriptsize] at (90:1.95) {Confidencialidade};
        \node[caixaambar,font=\scriptsize] at (210:2.0) {Integridade};
        \node[caixaescura,font=\scriptsize] at (330:2.0) {Disponibilidade};
        \node[font=\scriptsize\bfseries,text=azulpetroleo] at (0,0) {CIA};
      \end{tikzpicture}
      \\[2pt]{\tiny Tríade CID (em inglês, CIA).}
    \end{column}
    \begin{column}{0.56\textwidth}
      \scriptsize
      \begin{tabularx}{\linewidth}{@{}l L@{}}
        \toprule
        \cabecalho{Propriedade} & \cabecalho{Pergunta} \\
        \midrule
        Confidencialidade & só quem pode vê os dados? \\
        Integridade & os dados não foram alterados indevidamente? \\
        Disponibilidade & o serviço funciona quando precisam dele? \\
        Autenticidade & quem diz ser é quem é? \\
        \bottomrule
      \end{tabularx}
      \begin{conceito}[Princípios de projeto]
        \textbf{Menor privilégio} (cada parte com o mínimo de acesso), \textbf{defesa em profundidade} (várias camadas), \textbf{seguro por padrão} e \textbf{nunca confiar no cliente}: o app pode ser modificado por quem o instala.
      \end{conceito}
    \end{column}
  \end{columns}
\end{frame}

\begin{frame}{Criptografia essencial}
  \footnotesize
  \begin{tabularx}{\textwidth}{@{}p{2.5cm} L L L@{}}
    \toprule
    & \cabecalho{Hash} & \cabecalho{Simétrica} & \cabecalho{Assimétrica} \\
    \midrule
    Ideia & resumo de tamanho fixo, impossível de reverter & a mesma chave cifra e decifra & par de chaves: pública e privada \\
    Exemplos & SHA-256; para senhas, Argon2 ou bcrypt & AES-GCM & RSA, curvas elípticas (ECDSA, X25519) \\
    Uso no app & integridade de arquivos, armazenar senhas no servidor & dados locais cifrados, conteúdo do canal TLS & assinatura do app, troca de chaves, \emph{passkeys} \\
    \bottomrule
  \end{tabularx}
  \vspace{4pt}
  \begin{columns}[T]
    \begin{column}{0.48\textwidth}
      \begin{pordentro}[Como o HTTPS junta tudo]
        O TLS usa criptografia assimétrica para autenticar o servidor (certificado) e combinar uma chave secreta; depois cifra o tráfego com criptografia simétrica, que é muito mais rápida.
      \end{pordentro}
    \end{column}
    \begin{column}{0.48\textwidth}
      \begin{erro}[Criptografia caseira]
        Inventar algoritmos, reutilizar vetores de inicialização ou guardar a chave junto com o dado cifrado. Use bibliotecas e o cofre de chaves do sistema (Keystore, Keychain).
      \end{erro}
    \end{column}
  \end{columns}
\end{frame}

\begin{frame}{Modelagem de ameaças com STRIDE}
  \begin{columns}[T]
    \begin{column}{0.62\textwidth}
      \scriptsize
      \begin{tabularx}{\linewidth}{@{}l L L@{}}
        \toprule
        \cabecalho{Ameaça} & \cabecalho{Exemplo no LabMonitor} & \cabecalho{Defesa} \\
        \midrule
        \textbf{S}poofing (falsificação) & alguém se passa por um ESP32 e envia leituras falsas & autenticar dispositivos (certificados, MQTT com senha) \\
        \textbf{T}ampering (adulteração) & alterar a reserva no tráfego & TLS, validação no servidor \\
        \textbf{R}epudiation (repúdio) & aluno nega ter cancelado a reserva & registros de auditoria \\
        \textbf{I}nformation disclosure & token salvo em texto puro & armazenamento seguro \\
        \textbf{D}enial of service & app em laço enviando milhares de pedidos & limites de taxa no servidor \\
        \textbf{E}levation of privilege & aluno vira administrador alterando o app & regras e papéis no servidor \\
        \bottomrule
      \end{tabularx}
    \end{column}
    \begin{column}{0.36\textwidth}
      \begin{definicao}[Modelagem de ameaças]
        Analisar sistematicamente o que pode dar errado, \textbf{antes} de construir: desenhar o fluxo de dados, listar ameaças, priorizar e definir defesas.
      \end{definicao}
      \begin{historia}
        STRIDE foi criado na Microsoft no fim dos anos 1990 e virou uma das técnicas mais ensinadas.
      \end{historia}
    \end{column}
  \end{columns}
\end{frame}

\begin{frame}{OWASP Mobile Top 10 e MASVS}
  \begin{columns}[T]
    \begin{column}{0.5\textwidth}
      \scriptsize
      \begin{tabularx}{\linewidth}{@{}l L@{}}
        \toprule
        \multicolumn{2}{@{}l}{\cabecalho{OWASP Mobile Top 10 (2024)}} \\
        \midrule
        M1 & uso inadequado de credenciais \\
        M2 & segurança insuficiente da cadeia de suprimentos \\
        M3 & autenticação e autorização inseguras \\
        M4 & validação insuficiente de entradas e saídas \\
        M5 & comunicação insegura \\
        M6 & controles de privacidade inadequados \\
        M7 & proteção insuficiente do binário \\
        M8 & configuração de segurança incorreta \\
        M9 & armazenamento de dados inseguro \\
        M10 & criptografia insuficiente \\
        \bottomrule
      \end{tabularx}
    \end{column}
    \begin{column}{0.48\textwidth}
      \begin{definicao}[OWASP]
        Fundação aberta que publica guias de segurança. O \textbf{Top 10} lista os riscos mais comuns; o \textbf{MASVS} define os requisitos de segurança para apps, e o \textbf{MASTG} ensina a testá-los.
      \end{definicao}
      \begin{conceito}[Grupos do MASVS]
        Armazenamento, criptografia, autenticação, rede, plataforma, código, resiliência e privacidade.
      \end{conceito}
      \begin{dica}
        Use o MASVS como \emph{checklist} de revisão antes de cada versão importante.
      \end{dica}
    \end{column}
  \end{columns}
\end{frame}

% =============================================================================
\section{Fundamentos de produto e negócio}

\begin{frame}{Do problema ao MVP}
  \begin{columns}[c]
    \begin{column}{0.4\textwidth}
      \centering
      \begin{tikzpicture}
        \node[circle,fill=azulpetroleo,text=white,font=\scriptsize\bfseries,minimum size=1.3cm] (c) at (90:1.3) {Construir};
        \node[circle,fill=ambar,text=azulpetroleo,font=\scriptsize\bfseries,minimum size=1.3cm] (m) at (-30:1.3) {Medir};
        \node[circle,fill=azulprofundo,text=white,font=\scriptsize\bfseries,minimum size=1.3cm] (a) at (210:1.3) {Aprender};
        \draw[setaambar] (c) to[bend left=25] node[rotulo,right,font=\tiny]{produto} (m);
        \draw[setaambar] (m) to[bend left=25] node[rotulo,below,font=\tiny]{dados} (a);
        \draw[setaambar] (a) to[bend left=25] node[rotulo,left,font=\tiny]{ideias} (c);
      \end{tikzpicture}
    \end{column}
    \begin{column}{0.58\textwidth}
      \begin{definicao}[MVP --- Produto Mínimo Viável]
        A menor versão do produto que permite \textbf{aprender} com usuários reais se a solução resolve o problema. Não é um produto mal feito: é um produto pequeno.
      \end{definicao}
      \begin{conceito}[Métricas AARRR]
        \textbf{A}quisição (chegam), \textbf{A}tivação (têm uma boa primeira experiência), \textbf{R}etenção (voltam), \textbf{R}eceita (pagam) e \textbf{R}ecomendação (indicam).
      \end{conceito}
      \begin{historia}
        O ciclo construir--medir--aprender foi popularizado por Eric Ries no livro \emph{A Startup Enxuta} (2011).
      \end{historia}
    \end{column}
  \end{columns}
\end{frame}

\begin{frame}{Modelos de receita e métricas do negócio}
  \begin{columns}[T]
    \begin{column}{0.46\textwidth}
      \scriptsize
      \begin{tabularx}{\linewidth}{@{}l L@{}}
        \toprule
        \cabecalho{Modelo} & \cabecalho{Como gera receita} \\
        \midrule
        Pago & paga-se para instalar \\
        Freemium & grátis, com recursos pagos \\
        Assinatura & cobrança recorrente \\
        Anúncios & exibições e cliques \\
        B2B / institucional & a organização paga pelo serviço \\
        Público / acadêmico & financiado por projeto ou órgão \\
        \bottomrule
      \end{tabularx}
    \end{column}
    \begin{column}{0.52\textwidth}
      \scriptsize
      \begin{tabularx}{\linewidth}{@{}l L@{}}
        \toprule
        \cabecalho{Métrica} & \cabecalho{Definição} \\
        \midrule
        ARPU & receita média por usuário ativo no período \\
        \emph{Churn} & \% de assinantes que cancelam por mês \\
        LTV & receita esperada de um usuário durante toda a relação; aproximação: ARPU mensal $\times$ margem $\div$ \emph{churn} mensal \\
        CAC & custo para adquirir um usuário (marketing $\div$ novos usuários) \\
        \bottomrule
      \end{tabularx}
      \begin{exemplo}[Conta rápida]
        ARPU de R\$\,10/mês, margem de 70\,\% e \emph{churn} de 5\,\%/mês: LTV $\approx 10 \times 0{,}7 \div 0{,}05 =$ R\$\,140. Gastar R\$\,200 para adquirir cada usuário dá prejuízo.
      \end{exemplo}
    \end{column}
  \end{columns}
\end{frame}


% =============================================================================
\section{IA no app Flutter}

\begin{frame}{Onde roda o modelo?}
  \begin{columns}[T]
    \begin{column}{0.5\textwidth}
      \scriptsize
      \begin{tabularx}{\linewidth}{@{}l L L@{}}
        \toprule
        & \cabecalho{No aparelho} & \cabecalho{Na nuvem} \\
        \midrule
        Latência & muito baixa & depende da rede \\
        Offline & sim & não \\
        Privacidade & dados ficam no aparelho & dados saem do aparelho \\
        Custo & nenhum por uso & por requisição (\emph{tokens}) \\
        Modelo & pequeno (KB a centenas de MB) & muito grande \\
        \bottomrule
      \end{tabularx}
    \end{column}
    \begin{column}{0.48\textwidth}
      \scriptsize
      \begin{tabularx}{\linewidth}{@{}l L@{}}
        \toprule
        \cabecalho{Pacote} & \cabecalho{O que oferece} \\
        \midrule
        \cod{google\_mlkit\_*} & modelos prontos: texto (OCR), códigos, rostos, rótulos, tradução \\
        \cod{tflite\_flutter} & roda o \textbf{seu} modelo \cod{.tflite} (LiteRT) \\
        \cod{firebase\_ai} & modelos Gemini na nuvem, sem expor chaves \\
        \bottomrule
      \end{tabularx}
    \end{column}
  \end{columns}
  \vspace{4pt}
  \begin{columns}[T]
    \begin{column}{0.48\textwidth}
      \begin{conceito}[LiteRT]
        Novo nome do \textbf{TensorFlow Lite}: executa modelos treinados em Python, com aceleração por GPU ou NPU quando disponível.
      \end{conceito}
    \end{column}
    \begin{column}{0.48\textwidth}
      \begin{dica}
        Comece pelo mais simples que resolve: regra estatística, depois modelo pronto, depois modelo próprio, e só então um LLM.
      \end{dica}
    \end{column}
  \end{columns}
\end{frame}

\begin{frame}[fragile]{Anomalias nas leituras e OCR}
  \begin{columns}[T]
    \begin{column}{0.54\textwidth}
\begin{codigo}[title=1) Linha de base: z-score em uma janela]{Dart}
Stream<double> anomalias(Stream<double> leituras,
    {int janela = 60, double limiar = 3}) async* {
  final buf = ListQueue<double>();
  await for (final x in leituras) {
    if (buf.length == janela) {
      final m = buf.reduce((a, b) => a + b) / janela;
      final v = buf.map((y) => (y - m) * (y - m))
          .reduce((a, b) => a + b) / janela;
      final dp = sqrt(v);
      if (dp > 0 && (x - m).abs() / dp > limiar) yield x;
      buf.removeFirst();
    }
    buf.addLast(x);
  }
}
\end{codigo}
      \begin{conceito}[z-score]
        Distância da leitura à média recente, em desvios-padrão. Acima de 3, é muito improvável.
      \end{conceito}
    \end{column}
    \begin{column}{0.44\textwidth}
\begin{codigo}[title=2) Modelo próprio (tflite\_flutter)]{Dart}
final it = await Interpreter.fromAsset(
    'assets/anomalia.tflite');
final saida = [List.filled(1, 0.0)];   // [1, 1]
it.run([janelaNormalizada], saida);    // [1, 60]
\end{codigo}
\begin{codigo}[title=3) OCR da placa do equipamento]{Dart}
final ocr = TextRecognizer(
    script: TextRecognitionScript.latin);
final r = await ocr.processImage(
    InputImage.fromFilePath(foto.path));
final patrimonio = RegExp(r'\d{6}')
    .firstMatch(r.text)?.group(0);
await ocr.close();
\end{codigo}
      \begin{dica}
        Compare qualquer modelo com a linha de base: às vezes ela basta.
      \end{dica}
    \end{column}
  \end{columns}
\end{frame}

\begin{frame}[fragile]{Gemini pelo Firebase AI Logic}
  \begin{columns}[T]
    \begin{column}{0.58\textwidth}
\begin{codigo}[title=data/services/assistente\_service.dart]{Dart}
import 'package:firebase_ai/firebase_ai.dart';

final _modelo = FirebaseAI.googleAI().generativeModel(
  model: nomeDoModelo,        // um modelo Gemini atual
  systemInstruction: Content.system(
    'Você é assistente técnico de um laboratório de '
    'eletroeletrônica. Responda em até 3 frases e diga '
    'quando os dados não forem suficientes.'),
);

Future<String> explicarAlerta(String leituras) async {
  final r = await _modelo.generateContent([
    Content.text('Leituras da bancada B3: $leituras. '
        'Por que o alerta de temperatura disparou?'),
  ]);
  return r.text ?? 'Sem resposta do assistente.';
}
\end{codigo}
    \end{column}
    \begin{column}{0.4\textwidth}
      \begin{conceito}[Por que pelo Firebase]
        O app não guarda chave de API: o Firebase intermedeia a chamada, e o \textbf{App Check} (Play Integrity no Android, App Attest no iOS) garante que só o seu app consome a cota.
      \end{conceito}
      \begin{atencao}
        Mostre a resposta como \textbf{sugestão}, identificada como gerada por IA, e trate falhas e limites de uso. Modelos e preços mudam: consulte a lista atual antes de escolher.
      \end{atencao}
    \end{column}
  \end{columns}
\end{frame}

% =============================================================================
\section{Segurança no app Flutter}

\begin{frame}{Revisão MASVS de um app Flutter}
  \footnotesize
  \begin{tabularx}{\textwidth}{@{}l L@{}}
    \toprule
    \cabecalho{Grupo} & \cabecalho{O que fazer no Flutter} \\
    \midrule
    \textbf{Armazenamento} & segredos em \cod{flutter\_secure\_storage}; nada sensível em \cod{shared\_preferences}, logs ou capturas \\
    \textbf{Criptografia} & bibliotecas padrão (\cod{cryptography}, cofres do sistema); nunca algoritmos próprios \\
    \textbf{Autenticação} & tokens curtos e renováveis; biometria só para desbloquear segredos locais \\
    \textbf{Rede} & HTTPS sempre; \emph{pinning} em apps sensíveis, com plano de troca do certificado \\
    \textbf{Plataforma} & validar \emph{deep links} e dados de QR codes; WebView sem JavaScript desnecessário \\
    \textbf{Código} & \cod{--obfuscate}; nenhum segredo no binário; dependências auditadas (\cod{dart pub outdated}) \\
    \textbf{Resiliência} & Firebase App Check ou verificação de integridade no servidor \\
    \textbf{Privacidade} & coletar o mínimo; declarações das lojas coerentes com o código \\
    \bottomrule
  \end{tabularx}
  \vspace{4pt}
  \begin{destaque}
    O cliente nunca é confiável: toda regra importante é verificada também no servidor.
  \end{destaque}
\end{frame}

\begin{frame}[fragile]{\emph{Pinning} de certificado}
  \begin{columns}[T]
    \begin{column}{0.56\textwidth}
\begin{codigo}[title=Só aceitar o certificado da API do LabMonitor]{Dart}
import 'dart:io';
import 'package:flutter/services.dart';
import 'package:http/io_client.dart';

Future<IOClient> clienteComPinning() async {
  final pem = await rootBundle.load('assets/certs/api.pem');
  final ctx = SecurityContext(withTrustedRoots: false)
    ..setTrustedCertificatesBytes(pem.buffer.asUint8List());
  return IOClient(HttpClient(context: ctx));
}
\end{codigo}
    \end{column}
    \begin{column}{0.42\textwidth}
      \begin{definicao}[\emph{Certificate pinning}]
        O app confia apenas em um certificado (ou chave) específico, e não em qualquer autoridade certificadora instalada no aparelho. Dificulta interceptar o tráfego com um certificado falso.
      \end{definicao}
      \begin{atencao}
        Certificados expiram. Fixe a chave da autoridade intermediária ou mantenha dois certificados válidos, e planeje a troca antes do vencimento, ou o app para de funcionar.
      \end{atencao}
    \end{column}
  \end{columns}
\end{frame}

% =============================================================================
\section{Monetização e outras superfícies}

\begin{frame}[fragile]{Compras no app}
  \begin{columns}[T]
    \begin{column}{0.56\textwidth}
\begin{codigo}[title=in\_app\_purchase (trechos)]{Dart}
final iap = InAppPurchase.instance;
iap.purchaseStream.listen((compras) async {
  for (final c in compras) {
    if (c.status == PurchaseStatus.purchased) {
      final ok = await api.validarCompra(
          c.verificationData.serverVerificationData);
      if (ok) liberarRecursoPremium();
    }
    if (c.pendingCompletePurchase) {
      await iap.completePurchase(c);
    }
  }
});

final r = await iap.queryProductDetails({'relatorios_pro'});
await iap.buyNonConsumable(purchaseParam:
    PurchaseParam(productDetails: r.productDetails.first));
\end{codigo}
    \end{column}
    \begin{column}{0.42\textwidth}
      \begin{conceito}[Regras das lojas]
        Bens digitais usados no app passam pelo sistema de compras da loja, que retém uma comissão; bens físicos seguem outras regras.
      \end{conceito}
      \begin{atencao}
        Valide cada compra \textbf{no servidor} com as APIs da Play e da App Store antes de liberar o conteúdo.
      \end{atencao}
      \begin{dica}
        Anúncios: \cod{google\_mobile\_ads}. Serviços como RevenueCat simplificam assinaturas nas duas lojas.
      \end{dica}
    \end{column}
  \end{columns}
\end{frame}

\begin{frame}{Além da tela do celular}
  \begin{columns}[T]
    \begin{column}{0.32\textwidth}
      \centering
      \begin{tikzpicture}
        \draw[petroleo,rounded corners=4pt,fill=petroleo!6] (0,0) rectangle (2.6,1.4);
        \node[font=\tiny\bfseries,text=petroleo] at (1.3,1.15) {LabMonitor};
        \node[font=\large\bfseries,text=ambar] at (1.3,0.65) {24,8\,°C};
        \node[font=\tiny,text=cinzamedio] at (1.3,0.25) {Bancada 3 $\cdot$ há 1 min};
      \end{tikzpicture}
      \begin{conceito}[Widgets da tela inicial]
        O pacote \cod{home\_widget} envia dados do Flutter; o visual é nativo (Glance no Android, WidgetKit no iOS).
      \end{conceito}
    \end{column}
    \begin{column}{0.32\textwidth}
      \centering
      \begin{tikzpicture}
        \draw[petroleo,line width=1.4pt,fill=petroleoesc] (0,0) circle (0.72);
        \node[font=\small\bfseries,text=white] at (0,0.1) {24,8°};
        \node[font=\tiny,text=ambarclaro] at (0,-0.3) {B3};
      \end{tikzpicture}
      \begin{conceito}[Relógios]
        Wear OS e Apple Watch costumam exigir apps ou extensões nativas; o app Flutter no celular envia os dados.
      \end{conceito}
    \end{column}
    \begin{column}{0.32\textwidth}
      \centering
      \begin{tikzpicture}
        \draw[petroleo,rounded corners=3pt,fill=petroleo!6] (0,0) rectangle (2.8,1.4);
        \fill[petroleo] (0,1.1) rectangle (2.8,1.4);
        \node[font=\tiny,white] at (1.4,1.25) {Android Auto / CarPlay};
        \node[font=\tiny,text=cinzatexto] at (1.4,0.6) {lista simples, voz};
      \end{tikzpicture}
      \begin{conceito}[Carro e TV]
        Usam \emph{templates} e regras de segurança próprias de cada sistema, com integração nativa.
      \end{conceito}
    \end{column}
  \end{columns}
  \begin{dica}
    A arquitetura em camadas permite reaproveitar repositórios e serviços em todas essas superfícies; só a apresentação muda.
  \end{dica}
\end{frame}

% =============================================================================
\section{Carreira e portfólio}

\begin{frame}{Do curso ao mercado}
  \begin{columns}[T]
    \begin{column}{0.48\textwidth}
      \begin{conceito}[Portfólio que convence]
        \begin{itemize}
          \item Dois ou três apps completos no GitHub, com README, capturas e arquitetura explicada.
          \item Pelo menos um app publicado (mesmo em teste aberto).
          \item Testes e CI visíveis no repositório.
          \item Um projeto com hardware (BLE, IoT) é um diferencial raro.
        \end{itemize}
      \end{conceito}
    \end{column}
    \begin{column}{0.48\textwidth}
      \begin{dica}[Continue aprendendo]
        \begin{itemize}
          \item Documentação oficial, \emph{Flutter samples} e o canal do Flutter no YouTube.
          \item Notas de versão do Flutter e do Dart a cada lançamento.
          \item Comunidades: GDG, Flutter Brasil, eventos locais.
          \item Contribuir com pacotes de código aberto no pub.dev.
        \end{itemize}
      \end{dica}
      \begin{atencao}[Assistentes de IA]
        Use-os para acelerar, mas entenda e revise cada linha: você responde pelo código que publica.
      \end{atencao}
    \end{column}
  \end{columns}
\end{frame}

% =============================================================================
\section{Projeto final}

\begin{frame}{Projeto final: requisitos}
  \begin{columns}[T]
    \begin{column}{0.5\textwidth}
      \begin{conceito}[Requisitos mínimos]
        \begin{enumerate}
          \item Quatro ou mais telas com \cod{go\_router} e Material 3.
          \item MVVM com repositórios, serviços e injeção de dependência.
          \item Banco local (Drift) e preferências.
          \item API REST própria ou BaaS (Módulos 16 e 17).
          \item Funcionamento offline nas telas principais.
          \item Um recurso do dispositivo (BLE, câmera, GPS\ldots).
          \item Testes de unidade e de widget; CI no GitHub.
          \item Acessibilidade, tema escuro e build de release.
        \end{enumerate}
      \end{conceito}
    \end{column}
    \begin{column}{0.48\textwidth}
      \begin{projeto}[Sugestões de tema]
        \begin{itemize}
          \item Concluir e publicar o \textbf{LabMonitor}.
          \item Empréstimo de equipamentos do laboratório com QR code.
          \item Medidor de consumo de energia com ESP32 e sensor de corrente.
          \item Acompanhamento de estágio ou TCC.
        \end{itemize}
      \end{projeto}
      \begin{dica}
        Resolva um problema real de alguém que você conhece e valide o MVP com essa pessoa.
      \end{dica}
    \end{column}
  \end{columns}
\end{frame}

\begin{frame}{Entregáveis, cronograma e avaliação}
  \begin{columns}[T]
    \begin{column}{0.48\textwidth}
      \scriptsize
      \begin{tabularx}{\linewidth}{@{}l L@{}}
        \toprule
        \cabecalho{Semana} & \cabecalho{Marco} \\
        \midrule
        1 & proposta: problema, público, protótipo no Figma, arquitetura e modelo de ameaças \\
        2 & navegação, UI e dados locais \\
        3 & backend, recurso do dispositivo e offline \\
        4 & testes, acessibilidade, release e apresentação \\
        \bottomrule
      \end{tabularx}
      \begin{conceito}[Entregáveis]
        Repositório no GitHub com README, build de release (AAB ou link de teste), vídeo de 3 minutos e apresentação de 10 minutos.
      \end{conceito}
    \end{column}
    \begin{column}{0.5\textwidth}
      \scriptsize
      \begin{tabularx}{\linewidth}{@{}L r@{}}
        \toprule
        \cabecalho{Critério de avaliação} & \cabecalho{Peso} \\
        \midrule
        Funcionalidades e requisitos atendidos & 25\,\% \\
        Arquitetura e qualidade do código & 20\,\% \\
        Experiência do usuário, UI e acessibilidade & 15\,\% \\
        Testes automatizados e CI & 15\,\% \\
        Recursos do dispositivo e integração & 10\,\% \\
        Documentação (README, diagramas) & 5\,\% \\
        Apresentação e demonstração & 10\,\% \\
        \midrule
        \textbf{Total} & \textbf{100\,\%} \\
        \bottomrule
      \end{tabularx}
      \begin{dica}
        Pesos sugeridos: ajuste ao plano de ensino.
      \end{dica}
    \end{column}
  \end{columns}
\end{frame}

% =============================================================================
\section{Fechamento do curso}

\begin{frame}{A jornada completa}
  \centering
  \begin{tikzpicture}[x=1.0cm,
      etapa/.style={minimum width=1.82cm,minimum height=11mm,font=\tiny,text width=1.62cm,inner sep=2pt,align=center}]
    \foreach \i/\t/\s in {0/Ecossistema/caixa,1/Ambiente/caixa,2/Dart/caixa,3/Engenharia e Git/caixa,
        4/Design e UX/caixa,5/Widgets/caixa,6/Assíncrono/caixa}{
      \node[etapa,\s] (a\i) at (\i*1.98,0) {\textbf{M0\i}\\ \t};}
    \foreach \i/\t/\s [count=\k from 7] in {0/Estado e navegação/caixa,1/Arquitetura/caixapetroleo,2/Banco de dados/caixapetroleo,
        3/Persistência/caixapetroleo,4/Rede e APIs/caixapetroleo,5/Firebase/caixapetroleo,6/Dispositivo e IoT/caixapetroleo}{
      \node[etapa,\s] (b\i) at (\i*1.98,-1.6) {\textbf{M\ifnum\k<10 0\fi\k}\\ \t};}
    \foreach \i/\t/\s [count=\k from 14] in {0/Testes/caixapetroleo,1/Publicação/caixapetroleo,2/Backend BaaS/caixaescura,
        3/Backend próprio/caixaescura,4/Engenharia avançada/caixaescura,5/Nativo e multi\-plataforma/caixaescura,6/Projeto final/caixaambar}{
      \node[etapa,\s] (c\i) at (\i*1.98,-3.2) {\textbf{M\k}\\ \t};}
    \foreach \i [evaluate=\i as \j using int(\i+1)] in {0,...,5}{
      \draw[seta] (a\i) -- (a\j); \draw[seta] (b\i) -- (b\j); \draw[seta] (c\i) -- (c\j);}
    \draw[seta,rounded corners] (a6.south) -- ++(0,-2mm) -| (b0.north);
    \draw[seta,rounded corners] (b6.south) -- ++(0,-2mm) -| (c0.north);
  \end{tikzpicture}
  \vspace{2pt}
  \begin{resumo}[O que você construiu]
    Um app Flutter completo e publicado, com arquitetura profissional, dados locais e remotos, backend, integração com hardware, testes e entrega contínua, rodando no celular, na web e no desktop.
  \end{resumo}
\end{frame}

\begin{frame}{Glossário}
  \begin{glossario}
    Aprendizado de máquina & técnica em que o sistema aprende padrões a partir de dados \\
    Inferência & uso de um modelo treinado para responder a entradas novas \\
    Quantização & redução da precisão dos pesos (ex.: 32 para 8 bits) para encolher o modelo \\
    LLM & modelo de linguagem que gera texto prevendo o próximo \emph{token} \\
    RAG & técnica que inclui documentos recuperados na pergunta ao LLM \\
    STRIDE & classificação de ameaças: falsificação, adulteração, repúdio, vazamento, negação de serviço, elevação de privilégio \\
    MASVS & padrão da OWASP com os requisitos de segurança de apps móveis \\
    MVP & menor versão do produto que permite aprender com usuários reais \\
    LTV / CAC & valor de um cliente ao longo da relação / custo para adquiri-lo \\
  \end{glossario}
\end{frame}

\begin{frame}{Questões de revisão}
  \begin{questoes}
    \item Diferencie treinamento de inferência. Onde cada um costuma acontecer num app?
    \item Por que precisão e revocação são mais úteis que a acurácia para avaliar um detector de falhas?
    \item O que é alucinação em LLMs e como o app deve lidar com ela?
    \item Por que a chave de API de um LLM não pode ficar dentro do app? Qual a alternativa?
    \item Aplique STRIDE a uma funcionalidade do seu projeto final e proponha uma defesa por ameaça.
    \item Por que a compra no app deve ser validada no servidor?
    \item Qual o risco do \emph{certificate pinning} e como mitigá-lo?
    \item Calcule o LTV com ARPU de R\$\,8/mês, margem de 60\,\% e \emph{churn} de 4\,\%/mês.
  \end{questoes}
\end{frame}

\begin{frame}{Exercícios finais}
  \begin{exercicio}[Exercício 1 --- IA]
    Gere 24\,h de leituras simuladas com picos artificiais e compare o detector por z-score com um modelo treinado em Python e convertido para LiteRT: medir precisão e revocação de cada um.
  \end{exercicio}
  \begin{exercicio}[Exercício 2 --- Segurança]
    Revise o LabMonitor com o \emph{checklist} MASVS e o STRIDE e registre ao menos três correções como \emph{issues} no GitHub.
  \end{exercicio}
  \begin{exercicio}[Exercício 3 --- Assistente]
    Implemente ``explique este alerta'' com o Firebase AI Logic, marcando a resposta como gerada por IA.
  \end{exercicio}
  \begin{projeto}[Projeto final]
    Escolha o tema, entregue a proposta na semana 1 e siga o cronograma de quatro semanas.
  \end{projeto}
\end{frame}

\begin{frame}{Para aprofundar}
  \small
  \begin{itemize}
    \item ML Kit para Flutter: \url{https://pub.dev/packages/google_ml_kit}
    \item LiteRT: \url{https://ai.google.dev/edge/litert}
    \item Firebase AI Logic: \url{https://firebase.google.com/docs/ai-logic}
    \item OWASP MASVS e MASTG: \url{https://mas.owasp.org}
    \item Compras no app: \url{https://pub.dev/packages/in_app_purchase}
    \item Segurança no Flutter: \url{https://docs.flutter.dev/security}
  \end{itemize}
\end{frame}

\begin{frame}{Respostas das questões de revisão}
  \begin{respostas}
    \item Treinamento ajusta os pesos com muitos dados (servidor com GPU); inferência usa o modelo pronto para uma entrada nova (no aparelho ou na nuvem, a cada uso).
    \item Falhas são raras: um detector que nunca alerta acerta quase sempre, mas é inútil. Precisão mede alarmes falsos; revocação mede falhas não detectadas.
    \item É uma resposta fluente e falsa. O app deve apresentá-la como sugestão, citar fontes quando houver (RAG) e pedir confirmação humana para decisões.
    \item O binário pode ser extraído e a chave, usada por terceiros às suas custas. A chamada deve passar por um backend ou serviço intermediário (Firebase AI Logic com App Check).
    \item Exemplo para a reserva: falsificação de usuário (autenticação), adulteração da reserva (TLS e validação no servidor), elevação de privilégio (papéis verificados no servidor).
    \item O app pode ser modificado para fingir a compra; só o servidor, consultando a loja, confirma que o pagamento existe.
    \item Se o certificado mudar, o app para de se conectar. Mitigação: fixar a chave da autoridade intermediária, manter um certificado reserva e planejar a troca.
    \item LTV $\approx 8 \times 0{,}6 \div 0{,}04 =$ R\$\,120.
  \end{respostas}
\end{frame}

\begin{frame}[plain,noframenumbering]
  \begin{tikzpicture}[remember picture,overlay]
    \fill[azulpetroleo] (current page.north west) rectangle (current page.south east);
    \draw[ambar,line width=3pt] ([xshift=1.2cm,yshift=1.6cm]current page.west) -- ++(3.5,0);
    \node[anchor=west,text width=0.8\paperwidth] at ([xshift=1.2cm,yshift=0.6cm]current page.west)
      {\color{white}\Huge\bfseries Obrigado!};
    \node[anchor=west,text width=0.72\paperwidth] at ([xshift=1.25cm,yshift=-0.9cm]current page.west)
      {\color{ambarclaro}\normalsize \autorcurso\\[2pt]
       \color{white!75}\footnotesize Fim do curso \nomecurso{} --- trilha Flutter};
  \end{tikzpicture}
\end{frame}

\end{document}
